% !PRINTED 118
\weberpage{118}
%<<CONT

\[
 y=\frac{-(u+v)\pm\sqrt{-3(u+v)^2+12uv}}{2}
  =\frac{-(u+v)\pm\sqrt{-3}(u-v)}{2}.
\]
We therefore put
\begin{equation}
\varepsilon=\frac{-1+\sqrt{-3}}{2},
\tag{11}
\end{equation}
from which
\begin{equation}
\varepsilon^3=1,\qquad \varepsilon^2+\varepsilon+1=0,
\qquad
\varepsilon^2=\frac{-1-\sqrt{-3}}{2}
\tag{12}
\end{equation}
follows. Thus one obtains the three roots of the cubic equation (6)
% [sic] Source-original mathematical cross-reference: the printed text says (6), although the cubic equation itself is (1).
\begin{equation}
\begin{array}{rcl}
\alpha&=&u+v,\\
\beta&=&\varepsilon u+\varepsilon^2v,\\
\gamma&=&\varepsilon^2u+\varepsilon v.
\end{array}
\tag{13}
\end{equation}
Here $\varepsilon$ is a third root of unity different from $1$, expressed algebraically without the trigonometric functions. If one substitutes $\varepsilon u$ or $\varepsilon^2u$ for $u$, and hence $\varepsilon^2v$ or $\varepsilon v$ for $v$, the three quantities $\alpha,\beta,\gamma$ are interchanged with one another.

The expression (9) for the root of a cubic equation is called \gesperrt{Cardano's formula}.

\begin{center}
{\large \S.\ 36.}

\vspace{0.5em}
{\large\gesperrt{Cayley's Expression for Cardano's Formula.}}
\end{center}

Every cube root has, as we have seen, three different values, which one obtains from one of them by multiplication by $1,\varepsilon,\varepsilon^2$. Thus each of the two quantities $u,v$, as they are defined by (8), \S\,35, also has three values, and the sum $u+v$, if one takes account only of the conditions (8), has \gesperrt{nine different} values. Of these, however, only three give solutions of the cubic equation, and only by adjoining the relation (4), \S\,35,
\begin{equation}
3uv=-a
\tag{1}
\end{equation}
are the usable values singled out from among the nine.
